%% Question bank chapter file - generated by bank_tools.py build (do not edit while the bank is still being built:
%% it is regenerated from the per-test data in Question Bank\data). Every question is one bankq block:
%%   \begin{bankq}[id=..., source={...}, paper=.., qno=.., type=..., subtopic={...}, answer={...}, answer_src=key|solved,
%%                 working={...}, difficulty=1-10, diff_parts={K S I C F}, diff_note={...},
%%                 group=<passage id or empty>, also_in={duplicate copies in other tests}]
%% difficulty = min(10, K+S+I+C+F-1): K concepts 1-3, S steps 1-4, I insight 0-2, C computation 0-2, F format 0-1
%%   ...question LaTeX, options, figures (\bfig / \bfigi -> img/)...
%%   \end{bankq}
%% Passages shared by several questions are bankpassage blocks; their questions carry group=<passage id>.
\documentclass[10pt]{article}
\usepackage[a4paper,margin=1.8cm]{geometry}
\usepackage[T1]{fontenc}
\usepackage{amsmath,amssymb,mathtools}
\usepackage{graphicx}
\usepackage[export]{adjustbox}
\usepackage{xcolor,needspace,keyval}
\usepackage{enumitem}
\usepackage{booktabs,array,tabularx}
\usepackage[hidelinks]{hyperref}
\setlength{\parindent}{0pt}
\setlength{\parskip}{3pt}
\makeatletter
\define@key{bq}{id}{\def\bq@id{#1}}\define@key{bq}{source}{\def\bq@src{#1}}\define@key{bq}{paper}{\def\bq@paper{#1}}
\define@key{bq}{qno}{\def\bq@qno{#1}}\define@key{bq}{type}{\def\bq@type{#1}}\define@key{bq}{subtopic}{\def\bq@sub{#1}}
\define@key{bq}{answer}{\def\bq@ans{#1}}\define@key{bq}{answer_src}{\def\bq@asrc{#1}}\define@key{bq}{working}{\def\bq@work{#1}}
\define@key{bq}{group}{\def\bq@group{#1}}\define@key{bq}{also_in}{\def\bq@also{#1}}
\define@key{bq}{difficulty}{\def\bq@diff{#1}}\define@key{bq}{diff_parts}{\def\bq@dparts{#1}}\define@key{bq}{diff_note}{\def\bq@dnote{#1}}
\newenvironment{bankq}[1][]{\def\bq@work{}\def\bq@also{}\def\bq@group{}\def\bq@diff{?}\def\bq@dparts{}\def\bq@dnote{}\setkeys{bq}{#1}%
  \par\needspace{4\baselineskip}\medskip\noindent{\small\textbf{\bq@id}\hfill\color{gray}[\bq@src, Paper \bq@paper, Q\bq@qno]
  \bq@type\ $\cdot$ \bq@sub\ $\cdot$ \textbf{D\bq@diff}}\par\nopagebreak\relax}
  {\par{\small\color{teal}Answer: \textbf{\bq@ans}\ifx\bq@work\@empty\else\ (solved: \bq@work)\fi
  \ifx\bq@also\@empty\else\par\color{gray}Also in: \bq@also\fi}\par}
\define@key{bp}{id}{\def\bp@id{#1}}\define@key{bp}{source}{\def\bp@src{#1}}\define@key{bp}{questions}{\def\bp@qs{#1}}
\newenvironment{bankpassage}[1][]{\setkeys{bp}{#1}\par\needspace{4\baselineskip}\bigskip\noindent
  {\small\textbf{Passage \bp@id}\hfill\color{gray}[\bp@src; questions \bp@qs]}\par\nopagebreak\relax\itshape}{\par}
\makeatother
\newcommand{\bfig}[2]{\par\smallskip{\centering\includegraphics[width=#1pt,max width=0.62\linewidth,max height=0.32\textheight]{img/#2}\par}\smallskip}
\newcommand{\bfigi}[2]{\includegraphics[width=#1pt,max width=\linewidth,max height=0.32\textheight,valign=c]{img/#2}}
\newcommand{\bmissing}[1]{\textit{[figure missing in the source conversion: \detokenize{#1}]}}
\newenvironment{qoptions}{\begin{itemize}[leftmargin=1.8em,itemsep=1pt,topsep=2pt,label={}]}{\end{itemize}}
\newcommand{\opt}[2]{\item[(#1)] #2}
\providecommand{\fwstart}{}\providecommand{\fwend}{}\providecommand{\qfullsplit}{}
\newcolumntype{L}[1]{>{\raggedright\arraybackslash}p{#1}}
\newcolumntype{C}[1]{>{\centering\arraybackslash}p{#1}}
\newcolumntype{R}[1]{>{\raggedleft\arraybackslash}p{#1}}
\providecommand{\sectionhdr}[1]{}\providecommand{\subsectionhdr}[1]{}


\begin{document}
{\Large\bfseries Matrices --- Question Bank}\par\smallskip
{\small Mathematics chapter M07. Questions from all converted test papers, grouped by source series. Answers: printed key where it matches the paper, otherwise solved (answer\_src=solved, with working). Exact repeats across papers appear once, with also\_in listing the other copies.}\par\medskip\hrule

\section*{Score-I papers}
\begin{bankq}[id=SI12-P1-MATH-S1-Q08, source={Score-I - Test 12 (09 Mar 2022)}, paper=1, qno=8, type=multi, subtopic={Operations}, answer={A,B,D}, answer_src=key, difficulty=4, diff_parts={K1 S2 I0 C1 F1}, diff_note={concepts: Cayley-Hamilton theorem | steps: compute A^2; verify the polynomial relation; A^-1 from it; invertibility of A^2 and A^3 from det A != 0}]
Let $A=\begin{bmatrix}1&2&2\\2&1&2\\2&2&1\end{bmatrix}$, then
\begin{qoptions}
\opt{A}{$A^2-4A-5I_3=0$}
\opt{B}{$A^{-1}=\dfrac15(A-4I_3)$}
\opt{C}{$A^3$ is not invertible}
\opt{D}{$A^2$ is invertible}
\end{qoptions}
\end{bankq}
\begin{bankq}[id=SI12-P1-MATH-S1-Q09, source={Score-I - Test 12 (09 Mar 2022)}, paper=1, qno=9, type=multi, subtopic={Linear systems}, answer={B,C,D}, answer_src=key, difficulty=3, diff_parts={K1 S2 I0 C0 F1}, diff_note={concepts: consistency of a linear system | steps: equations 1 and 3 need mu = 6 for solutions; case lambda = 3; case lambda = 2; case lambda = 4 and 5}]
System of equation\\ $x+3y+2z=6$\\ $x+\lambda y+2z=7$\\ $x+3y+2z=\mu$ has
\begin{qoptions}
\opt{A}{unique solution if $\lambda=2$, $\mu\ne6$}
\opt{B}{infinitely many solution if $\lambda=4$, $\mu=6$}
\opt{C}{no solution if $\lambda=5$, $\mu=7$}
\opt{D}{no solution if $\lambda=3$, $\mu=5$}
\end{qoptions}
\end{bankq}
\begin{bankq}[id=SI12-P2-MATH-S2-Q06, source={Score-I - Test 12 (09 Mar 2022)}, paper=2, qno=6, type=single, subtopic={Operations}, answer={1.00}, answer_src=key, difficulty=4, diff_parts={K2 S1 I2 C0 F0}, diff_note={concepts: trace of powers, nilpotent matrices | steps: multiply AC - CA = C^2 by C^k and take traces so all tr(C^k) = 0; C is nilpotent so det C = 0 (one value) | I: trace argument}]
Let $A$ and $C$ are square matrices of same order, $\det(A)\ne0$ and $AC-CA=C^2$, then the number of values of $\det(C)$ is
\end{bankq}
\end{document}
